An increase in the output price can raise the income left for workers and owners,
but the amount also depends on the price of the input the firm buys. Consider a
competitive firm with fixed labor $L$ and production $q=f(L,x)$. The firm chooses
the input $x$ after observing the output price $p$ and input price $r$. Its
restricted profit is
\[
 G(p,r;L)=\max_{x\geq0}\{p f(L,x)-rx\}.
\]
Suppose the optimum exists, is interior, and the production and value functions
are differentiable at the prices considered. Concavity in the purchased input
provides a sufficient condition for the first-order condition to characterize
the optimum. With labor fixed, that condition is $p f_x(L,x)=r$.

To see what input adjustment contributes, differentiate the optimized payments:
\[
 dG=q\,dp-x\,dr+[p f_x-r]\,dx=q\,dp-x\,dr.
\]
An extra unit of the input raises sales by exactly its purchase cost at the
initial optimum. The induced quantity change therefore contributes no separate
first-order gain. The output price rise values initial production; an input
price rise reduces income by its value on initial input purchases. Neither
conclusion assumes a particular functional form for production.

When the input price is unchanged, this envelope also relates income to sales.
Let $S=pq$ and consider a change in $p$ with $r$ and $L$ fixed. Provided $dS$ is
nonzero, define $\kappa=d\log p/d\log S$. Then $dG=\kappa\,dS$. This marginal
price share depends on the quantity response along that change. It is not the
average share $G/S$, which describes how initial sales are divided rather than
how additional sales arise. If $r$ changes too, the input-price term remains in
the envelope and that one-dimensional price mapping no longer gives the whole
income change. Allocating the change between workers and owners requires a
further account of their claims.
